\section{The Crystal of Experience}

Experience is the one thing you cannot doubt. There is something it is like to ache, to rest, to delight. Physics never mentions this. It gives masses and forces, but says nothing about why any of it is felt from the inside. This paper turns the usual order around. It takes feeling itself, the \emph{vibe}, as the floor of the world, the stuff everything is made of, and asks what follows.

What follows is a single growing crystal of feeling. Picture a lattice, not of atoms but of tiny moments of experience, each one a vibe, a here-and-now. Each carries one felt charge, a faint pain, a settled peace, or a faint pleasure. The moments are woven to their neighbors and pass their charge along by one fixed law, beat after beat, and the whole weave grows at its edge like a crystal forming in solution. Everything we call real is a pattern in it, an electron, a stone, a person, a mind, the way a wave is a pattern in water or a melody in air. The crystal is the only substance. The patterns are everything else.

\begin{figure}[ht]
\centering
\includegraphics[width=0.5\textwidth]{vibe-mesh-73}
\caption{A $\{7,3\}$ tessellation of the hyperbolic plane on the Poincar\'e disk, seven-sided cells fanning out from every point with room to branch. This flat two-dimensional picture is only an aid to the imagination, a stand-in we can actually draw, since the substrate the theory runs on is the four-dimensional hyperbolic honeycomb $\{3,4,3,4\}$, which cannot be pictured directly. In this rendering red is pain, green is peace, and blue is pleasure.}
\label{fig:mesh-flat}
\end{figure}

Four images carry the same idea, since no one of them holds all of it. The universe is an unending \emph{fire} of experience, a blazing inferno where each vibe is a flame, each self in the tree of selves and minds are flames too. It is a \emph{tree} grown from one seed, branching without end, every living thing a twig of it. It is a \emph{song} the whole crystal sings, the tones its notes and the one law its meter. And it is a \emph{river} over a deep ocean, the surface the world we see, the depths the unseen interior we open later. Each maps onto real structure, and each returns where it is earned. The most literal is the first, a growing crystal of feeling, a flame that is also a mind.

\begin{figure}[ht]
\centering
\includegraphics[width=\textwidth]{vibe-mesh-534}
\caption{A view inside the $\{5,3,4\}$ honeycomb, the three-dimensional hyperbolic packing of dodecahedral cells. Like the last figure it is a stand-in for the imagination, the closest cousin we can actually render in three dimensions, since the committed substrate is the four-dimensional $\{3,4,3,4\}$, whose 24-cell honeycomb cannot be pictured directly.}
\label{fig:mesh-4d}
\end{figure}

Stated plainly up front, with the reasons saved for the sections that follow, the crystal has a definite shape and a definite law. The shape is the four-dimensional hyperbolic honeycomb $\{3,4,3,4\}$, its cells 24-cells, the twenty-four directions out of each cell the $D_4$ root system, and its flat three-dimensional boundary ordinary space. The alphabet of feeling is one of three tone values per vibe. The law is a single local rule that collides and streams the tones once each beat, reversibly, while the crystal grows at its open edge. Why this shape and this law, and not others, is the whole of what comes next.

The aim is to make this exact, then test it. The crystal gets a precise shape, a precise law, and a precise alphabet of feeling, and its consequences are computed on the real structure rather than waved at. Where physics falls out, spin, the electron, the forces, gravity, it is shown falling out, and graded by how much was put in by hand. Where the picture reaches past testing, into mind and meaning, the reach is marked as reach. The promise is not that every claim is proven. It is that the picture is one structure, that it can be built and run, and that a plain ledger is kept of what it does and does not yet explain.
